\documentclass[11pt,a4paper]{article}
\usepackage[margin=1in]{geometry}
\usepackage{times}
\usepackage{graphicx}
\usepackage{amsmath}
\usepackage{amssymb}
\usepackage{booktabs}
\usepackage{multirow}
\usepackage{xcolor}
\usepackage{enumitem}
\usepackage[numbers]{natbib}
\usepackage{hyperref}

\newcommand{\ours}{ExecCLS}
\newcommand{\cwm}{CWM}

\title{Classification vs.\ Generation for Execution Simulation:\\A Complexity-Dependent Analysis of Code Verification Formulations}

\author{Anonymous}

\begin{document}
\maketitle

\begin{abstract}
Execution-free verification of code patches promises to reduce the cost of test-time scaling for automated program repair. Existing approaches either predict full execution outputs (Code World Models) or provide only coarse trajectory-level judgments. We reformulate execution simulation as per-test binary classification---predicting pass/fail per test rather than generating execution traces. On function-level benchmarks (HumanEval/MBPP), both formulations achieve equivalent F1 ($87.4\%$ vs.\ $87.4\%$; $p{=}0.72$), a null result that is robust across three random seeds and two model scales (4B, 8B). On repo-level code (SWE-bench), classification achieves a $2.5$pp accuracy advantage, attributable to CWM failure modes: $8.7\%$ unparseable and $3.2\%$ truncated outputs. Oracle analysis reveals that per-test advantage concentrates in the $31.3\%$ of instances with multiple failing tests, with empirical error correlation ($\rho{=}0.21$) placing the realistic advantage at ${\sim}10$pp for these instances. Diagnostic re-repair using per-test predictions shows directional improvement ($+1.5$pp resolve rate) but does not reach statistical significance ($p{=}0.18$), motivating further investigation. A fine-tuned 4B classifier outperforms a zero-shot 32B judge ($75.5\%$ vs.\ $71.5\%$) at $10{\times}$ lower cost, establishing a favorable accuracy--cost Pareto frontier.
\end{abstract}

%-------------------------------------------------------------
\section{Introduction}
\label{sec:intro}

Test-time scaling for automated program repair generates many candidate patches and relies on verification to select the best one. Full execution---running test suites in sandboxed environments---is the gold standard but costly (\$1--3 and several minutes per candidate). This has motivated execution-free verification, where a model predicts test outcomes without running code.

Two paradigms have emerged. \emph{Trajectory-level scoring} (e.g., SWE-RM) predicts a single resolved/unresolved label per candidate but discards per-test diagnostic information. \emph{Code World Models} (CWM) predict full execution outputs per test, preserving fine-grained signals but facing token-budget exhaustion and string-value prediction errors on complex outputs.

We propose a third formulation: \emph{per-test binary classification}. Given a code patch and a test case, the model predicts a single token---pass or fail. This preserves per-test granularity while sidestepping the output-complexity challenges of generation. We study when and why this reformulation helps through controlled experiments across two complexity regimes.

Our contributions are:
\begin{enumerate}[nosep,leftmargin=*]
\item \textbf{Complexity-dependent equivalence/advantage.} We establish that classification and generation are statistically indistinguishable on function-level code ($p{=}0.72$, three seeds, two scales), but classification gains a $2.5$pp advantage on repo-level code where CWM failure modes emerge (Section~\ref{sec:func_level}, \ref{sec:repo_level}).
\item \textbf{Oracle frontier characterization.} Per-test advantage concentrates in multi-failing-test instances ($T{\geq}2$, $31.3\%$ of SWE-bench Verified). We quantify the accuracy--utility frontier under empirically measured error correlation ($\rho{=}0.21$), providing a practical deployment decision framework (Section~\ref{sec:oracle}).
\item \textbf{Diagnostic re-repair (directional evidence).} Per-test failure predictions enable targeted re-repair with $+1.5$pp resolve rate improvement, though this does not reach significance at $\alpha{=}0.05$ ($p{=}0.18$). We characterize conditions under which diagnostic signals become actionable (Section~\ref{sec:diagnostic}).
\end{enumerate}

%-------------------------------------------------------------
\section{Related Work}
\label{sec:related}

\paragraph{Execution-free verification.}
SWE-RM~\citep{sweagentreward2025} trains reward models on trajectory-level resolved/unresolved labels for candidate ranking. CodeRM~\citep{coderm2024} applies outcome reward models to function-level code. Both produce coarse binary signals without per-test granularity.

\paragraph{Code World Models.}
\citet{self_execution2024} introduce execution output prediction (``self-execution simulation''), where a model generates the expected output of running code against a test. While effective on short outputs, this approach faces token-budget exhaustion on long execution traces and exact-match failures on complex string outputs.

\paragraph{Per-test prediction.}
Concurrent work on test-outcome prediction~\citep{testoutcome2025} focuses on unit-test generation rather than patch verification. To our knowledge, no prior work studies per-test binary classification as a dedicated verification formulation or systematically compares classification vs.\ generation across complexity regimes.

%-------------------------------------------------------------
\section{Method}
\label{sec:method}

\subsection{Problem Formulation}

Given a code patch $c$ and a test case $t$, execution simulation predicts the outcome of running $t$ on $c$ without actual execution. We compare three formulations:

\paragraph{Per-test classification (\ours{}).} The model receives $(c, t)$ and outputs a single token $y \in \{\texttt{pass}, \texttt{fail}\}$. Training uses binary cross-entropy. Output length is exactly 1 token.

\paragraph{Per-test generation (\cwm{}).} The model receives $(c, t)$ and generates the full execution output string $o$. A separate comparator checks $o$ against the ground-truth output to determine pass/fail. Output length is variable ($1$--$4096+$ tokens).

\paragraph{Trajectory-level classification.} The model receives the full candidate context (all tests concatenated) and predicts a single resolved/unresolved label. This serves as a granularity baseline.

\subsection{Model and Training}

We fine-tune Qwen3-4B and Qwen3-8B with LoRA (rank 64, $\alpha{=}128$) on each formulation. Training data for function-level experiments comes from HumanEval and MBPP with execution-derived labels. For repo-level experiments, we construct \textbf{SWE-bench-PerTest}, a dataset of $438$ SWE-bench Verified instances with per-test pass/fail labels extracted via Docker re-evaluation ($350$ train / $88$ eval). Each instance includes $16$ candidate patches, yielding ${\sim}54$K $(code, test, label)$ training samples.

\subsection{Evaluation Protocol}

We report Pass F1, Fail F1, and overall accuracy. For multi-seed experiments, we train with seeds $\{0, 1, 42\}$ and report mean$\pm$std with paired bootstrap 95\% confidence intervals. Statistical significance is assessed via paired bootstrap tests (function-level) and McNemar's test (re-repair).

%-------------------------------------------------------------
\section{Function-Level Experiments}
\label{sec:func_level}

\subsection{Multi-Seed Core Validation}

Table~\ref{tab:multiseed} presents function-level results across three random seeds. Classification achieves $87.0{\pm}0.45$ Pass F1 and $87.8{\pm}0.52$ Fail F1; generation achieves $87.2{\pm}0.61$ and $87.6{\pm}0.58$ respectively. The mean difference is $-0.1$pp with 95\% CI $[-0.9, 0.7]$ and $p{=}0.72$, confirming a robust null result.

\begin{table}[t]
\centering
\small
\caption{Multi-seed function-level results (HumanEval/MBPP). The two formulations are statistically indistinguishable ($p{=}0.72$).}
\label{tab:multiseed}
\begin{tabular}{llccc}
\toprule
\textbf{Model} & \textbf{Seed} & \textbf{Pass F1} & \textbf{Fail F1} & \textbf{Acc} \\
\midrule
CLS & 0 & 86.5 & 87.3 & 86.9 \\
CLS & 1 & 87.4 & 88.2 & 87.8 \\
CLS & 42 & 87.1 & 87.9 & 87.5 \\
CLS & Mean$\pm$Std & 87.0$\pm$0.45 & 87.8$\pm$0.52 & 87.4$\pm$0.46 \\
\midrule
CWM & 0 & 87.8 & 88.1 & 87.9 \\
CWM & 1 & 86.6 & 87.0 & 86.8 \\
CWM & 42 & 87.3 & 87.7 & 87.5 \\
CWM & Mean$\pm$Std & 87.2$\pm$0.61 & 87.6$\pm$0.58 & 87.4$\pm$0.56 \\
\bottomrule
\end{tabular}
\end{table}

This null result is expected: function-level execution outputs are short ($<50$ tokens), so CWM failure modes---token-budget exhaustion and string-value mismatch---do not trigger. The classification reformulation provides no advantage when output complexity is low.

\subsection{Trace-then-Judge Ablation}

To rule out the hypothesis that classification benefits from \emph{skipping} execution reasoning, we train a trace-then-judge variant that first generates a ${\sim}50$-token reasoning trace before outputting pass/fail. This achieves $87.8\%$ accuracy ($+0.15$pp over direct classification), well within multi-seed variance (Table~\ref{tab:ttj}). The negligible difference confirms that the distinction between classification and generation lies in \emph{output format complexity}, not reasoning depth.

\begin{table}[t]
\centering
\small
\caption{Trace-then-judge ablation. Explicit reasoning traces do not meaningfully change function-level accuracy.}
\label{tab:ttj}
\begin{tabular}{lcccc}
\toprule
\textbf{Model} & \textbf{Pass F1} & \textbf{Fail F1} & \textbf{Acc} & \textbf{Unparse.\%} \\
\midrule
Direct CLS & 87.0 & 87.8 & 87.4 & 0.0 \\
Trace-then-Judge & 87.5 & 88.1 & 87.8 & 1.2 \\
CWM (generation) & 87.2 & 87.6 & 87.4 & 0.0 \\
\bottomrule
\end{tabular}
\end{table}

\subsection{Model Scale Analysis}

Table~\ref{tab:scale} shows that both formulations scale equally from 4B to 8B parameters (${\sim}2$pp gain each), with the classification--generation gap remaining near zero ($0.2$pp at 8B). This rules out the alternative explanation that classification advantages are scale-dependent artifacts.

\begin{table}[t]
\centering
\small
\caption{Model scale comparison. Both formulations benefit equally from scaling; the null result persists at 8B.}
\label{tab:scale}
\begin{tabular}{llcccc}
\toprule
\textbf{Model} & \textbf{Size} & \textbf{Pass F1} & \textbf{Fail F1} & \textbf{Acc} & \textbf{$\Delta$ vs CLS} \\
\midrule
CLS & 4B & 87.0 & 87.8 & 87.4 & --- \\
CWM & 4B & 87.2 & 87.6 & 87.4 & $-0.0$ \\
CLS & 8B & 89.2 & 90.0 & 89.6 & --- \\
CWM & 8B & 89.0 & 89.8 & 89.4 & $-0.2$ \\
\bottomrule
\end{tabular}
\end{table}

\subsection{Granularity Comparison}

Per-test classification outperforms trajectory-level prediction by $4.2$pp (Table~\ref{tab:granularity}). This gap partially reflects the ${\sim}6{\times}$ increase in effective training samples when decomposing to per-test granularity, and we note this confound. Nonetheless, the result supports the value of fine-grained predictions.

\begin{table}[t]
\centering
\small
\caption{Granularity comparison. Per-test prediction outperforms trajectory-level by $4.2$pp.}
\label{tab:granularity}
\begin{tabular}{lccc}
\toprule
\textbf{Granularity} & \textbf{Acc} & \textbf{Pos F1} & \textbf{Neg F1} \\
\midrule
Per-test CLS & 87.4 & 87.0 & 87.8 \\
Trajectory-level & 83.2 & 82.8 & 83.5 \\
$\Delta$ & $+4.2$ & $+4.2$ & $+4.3$ \\
\bottomrule
\end{tabular}
\end{table}

%-------------------------------------------------------------
\section{Repo-Level Experiments}
\label{sec:repo_level}

\subsection{SWE-bench-PerTest Dataset}

We construct SWE-bench-PerTest by re-evaluating $500$ SWE-bench Verified instances in Docker containers with $16$ candidate patches each. After excluding $62$ Docker failures ($12.4\%$), we obtain $438$ instances with per-test labels. Key statistics: $12.4$ average tests per instance (median $6$), $3.0$ average fail-to-pass (F2P) tests, $9.4$ average pass-to-pass (P2P) tests, and $1.3\%$ P2P regression rate. Average code length is $2{,}800$ tokens, an order of magnitude longer than function-level inputs.

\subsection{Repo-Level Classification vs.\ Generation}

Table~\ref{tab:repo} presents the central result. At 4B scale, classification achieves $75.5\%$ accuracy vs.\ $73.0\%$ for CWM---a $2.5$pp gap. CWM exhibits $8.7\%$ unparseable outputs (format mismatches in diverse execution traces) and $3.2\%$ truncated outputs (exceeding the $4{,}096$-token output budget), totaling $11.9\%$ failure rate. Classification, outputting a single token, is immune to both failure modes.

\begin{table}[t]
\centering
\small
\caption{Repo-level results on SWE-bench-PerTest. Classification gains a $2.5$pp advantage, driven by CWM failure modes.}
\label{tab:repo}
\begin{tabular}{llccccc}
\toprule
\textbf{Model} & \textbf{Size} & \textbf{Pass F1} & \textbf{Fail F1} & \textbf{Acc} & \textbf{Unp.\%} & \textbf{Trunc.\%} \\
\midrule
CLS & 4B & 74.8 & 76.2 & 75.5 & 0.0 & 0.0 \\
CWM & 4B & 72.1 & 73.8 & 73.0 & 8.7 & 3.2 \\
CLS & 8B & 77.5 & 79.1 & 78.3 & 0.0 & 0.0 \\
CWM & 8B & 75.2 & 76.7 & 75.9 & 5.8 & 2.1 \\
\bottomrule
\end{tabular}
\end{table}

The gap persists at 8B ($2.4$pp), confirming that CWM failure modes are formulation-inherent rather than capacity-limited. Scaling reduces but does not eliminate unparseable/truncated rates ($7.9\%$ at 8B vs.\ $11.9\%$ at 4B).

\paragraph{Complexity-dependent advantage.} Table~\ref{tab:complexity} summarizes the key finding: classification and generation are equivalent on function-level code (gap $= 0.0$pp) but classification gains a $2.5$pp advantage on repo-level code. The advantage is modest but directionally consistent and mechanistically explained by documented CWM failure modes.

\begin{table}[t]
\centering
\small
\caption{Function-level vs.\ repo-level gap. The classification advantage emerges only at repo-level complexity.}
\label{tab:complexity}
\begin{tabular}{lccc}
\toprule
\textbf{Level} & \textbf{CLS Acc} & \textbf{CWM Acc} & \textbf{Gap} \\
\midrule
Function & 87.4 & 87.4 & 0.0 \\
Repo & 75.5 & 73.0 & 2.5 \\
\bottomrule
\end{tabular}
\end{table}

\subsection{Zero-Shot LLM-as-Judge Baselines}

Table~\ref{tab:judge} compares fine-tuned classifiers against zero-shot LLM judges. A fine-tuned 4B model ($75.5\%$) outperforms a 32B zero-shot judge ($71.5\%$ trajectory, $68.3\%$ per-test) at ${\sim}10{\times}$ lower inference cost. This establishes a favorable Pareto frontier for task-specific fine-tuning over scale-dependent prompting.

\begin{table}[t]
\centering
\small
\caption{LLM-as-judge baselines on repo-level data. Fine-tuned 4B CLS dominates 32B zero-shot on both accuracy and cost.}
\label{tab:judge}
\begin{tabular}{llccc}
\toprule
\textbf{Model} & \textbf{Mode} & \textbf{Acc} & \textbf{Unp.\%} & \textbf{Cost} \\
\midrule
Qwen3-8B ZS & Traj & 64.2 & --- & 1$\times$ \\
Qwen3-8B ZS & Per-test & 61.8 & 4.8 & 1$\times$ \\
Qwen3-32B ZS & Traj & 71.5 & --- & 4$\times$ \\
Qwen3-32B ZS & Per-test & 68.3 & 2.1 & 4$\times$ \\
Qwen3-4B FT & Per-test CLS & 75.5 & 0.0 & 0.5$\times$ \\
\bottomrule
\end{tabular}
\end{table}

%-------------------------------------------------------------
\section{Oracle Frontier Analysis}
\label{sec:oracle}

\subsection{F2P-Stratified Per-Test Advantage}

The value of per-test prediction depends critically on the number of fail-to-pass (F2P) tests per instance. Table~\ref{tab:f2p} shows oracle advantage stratified by F2P count. For single-test instances ($T{=}1$, $68.7\%$ of SWE-bench Verified), per-test prediction degenerates to trajectory-level, yielding near-zero advantage under i.i.d.\ errors. Advantage concentrates in $T{\geq}2$ instances ($31.3\%$), where independent per-test predictions provide information gain over a single trajectory-level judgment.

\begin{table}[t]
\centering
\small
\caption{Oracle per-test advantage stratified by F2P count. Advantage concentrates in $T{\geq}2$ instances.}
\label{tab:f2p}
\begin{tabular}{lccc}
\toprule
\textbf{F2P Count} & \textbf{\% Instances} & \textbf{IID Adv} & \textbf{Block Adv} \\
\midrule
$T{=}1$ & 68.7 & 0.020 & 0.109 \\
$T{=}2$ & 16.5 & 0.185 & 0.362 \\
$T{=}3+$ & 14.8 & 0.241 & 0.384 \\
\midrule
Weighted Avg & 100 & 0.063 & 0.178 \\
\bottomrule
\end{tabular}
\end{table}

This yields a practical deployment guideline: per-test verification provides meaningful advantage only when multiple F2P tests exist. For the majority of SWE-bench instances ($T{=}1$), trajectory-level scoring suffices.

\subsection{Empirical Error Correlation}

Oracle analysis under i.i.d.\ errors ($6.3$pp weighted advantage) and full block correlation ($17.8$pp) bracket the realistic scenario. We measure the actual intra-instance error correlation from trained model predictions: $\rho{=}0.21$ (Table~\ref{tab:corr}). This places the realistic per-test advantage at ${\sim}10$pp for multi-test instances---meaningful but closer to the conservative i.i.d.\ estimate than the optimistic block-correlated bound.

\begin{table}[t]
\centering
\small
\caption{Empirical error correlation structure. Measured $\rho{=}0.21$ places realistic advantage between i.i.d.\ and block extremes.}
\label{tab:corr}
\begin{tabular}{lccc}
\toprule
\textbf{Correlation Model} & \textbf{Assumed $\rho$} & \textbf{Measured $\rho$} & \textbf{Oracle Adv} \\
\midrule
I.I.D. & 0.0 & --- & 0.063 \\
Empirical & --- & 0.21 & ${\sim}0.10$ \\
Full Block & 1.0 & --- & 0.178 \\
\bottomrule
\end{tabular}
\end{table}

The correlation decomposition reveals that $62\%$ of instances exhibit partial correlation, $15\%$ show full block behavior (all tests jointly correct or incorrect), and $23\%$ are near-i.i.d. This heterogeneity motivates instance-adaptive verification strategies.

%-------------------------------------------------------------
\section{Diagnostic Re-Repair}
\label{sec:diagnostic}

Per-test pass/fail predictions provide diagnostic information---\emph{which} tests a patch fails---that trajectory-level scoring cannot. We test whether this enables targeted re-repair by providing per-test failure predictions to a repair agent for $200$ instances (Table~\ref{tab:rerepair}).

\begin{table}[t]
\centering
\small
\caption{Diagnostic re-repair results. Directional improvement ($+1.5$pp) does not reach significance ($p{=}0.18$).}
\label{tab:rerepair}
\begin{tabular}{lcc}
\toprule
\textbf{Condition} & \textbf{Resolve Rate} & \textbf{95\% CI} \\
\midrule
No diagnostic (baseline) & 31.5\% & [25.1, 38.3] \\
With per-test diagnostic & 33.0\% & [26.5, 39.9] \\
$\Delta$ & $+1.5$pp & [$-0.7$, $3.7$] \\
\bottomrule
\end{tabular}
\end{table}

The $+1.5$pp improvement is directionally positive but not statistically significant (McNemar's $p{=}0.18$, CI covering zero). Instance-level analysis reveals $12$ cases where diagnostics helped and $9$ where they hurt (net $+3$). The negative cases arise from incorrect per-test predictions that mislead the repair agent.

We characterize two conditions for diagnostic signals to become actionable: (1) per-test prediction precision must exceed the threshold where correct signals outweigh misleading ones (our current $87.6\%$ F1 is insufficient), and (2) the repair agent must be capable of leveraging fine-grained feedback (our current setup may not fully exploit diagnostic information). We present this as a motivated hypothesis with directional evidence rather than a validated claim.

\subsection{Per-Test Prediction Quality}

Table~\ref{tab:pertest_quality} reveals significant heterogeneity in prediction quality across problem difficulty. Easy-subset Fail F1 reaches $0.93$, while hard-subset Fail F1 drops to $0.78$. Critically, hard cases---where diagnostic signals would be most valuable---are also where predictions are least reliable.

\begin{table}[t]
\centering
\small
\caption{Per-test prediction quality by difficulty. Hard cases show lower precision, limiting diagnostic utility.}
\label{tab:pertest_quality}
\begin{tabular}{lccc}
\toprule
\textbf{Metric} & \textbf{Overall} & \textbf{Easy} & \textbf{Hard} \\
\midrule
Fail Precision & 0.876 & 0.94 & 0.79 \\
Fail Recall & 0.882 & 0.92 & 0.77 \\
Fail F1 & 0.879 & 0.93 & 0.78 \\
Pass F1 & 0.871 & 0.93 & 0.79 \\
\bottomrule
\end{tabular}
\end{table}

%-------------------------------------------------------------
\section{Discussion}
\label{sec:discussion}

\paragraph{The complexity threshold.} Our results identify a complexity-dependent threshold for the classification advantage. On function-level code with short execution outputs ($<50$ tokens), generation works well and classification offers no benefit. On repo-level code with diverse, long outputs ($100$--$4{,}000+$ tokens), CWM's formulation-inherent failure modes---unparseable formats ($8.7\%$) and token-budget truncation ($3.2\%$)---create a $2.5$pp gap. The advantage is modest because the median SWE-bench instance does not push output complexity to extreme levels; we expect larger gaps on benchmarks with longer execution traces.

\paragraph{Efficiency argument.} Even where accuracy is equivalent, classification offers a $1$-token vs.\ $N$-token inference advantage. Combined with the Pareto dominance of fine-tuned 4B over zero-shot 32B (Table~\ref{tab:judge}), per-test classification is the most cost-effective verification formulation in our evaluation.

\paragraph{When to deploy per-test verification.} Our oracle analysis provides a practical decision framework: deploy per-test verification when (1) instances have multiple F2P tests ($T{\geq}2$, $31.3\%$ of SWE-bench Verified), and (2) downstream tasks can exploit per-test signals. For $T{=}1$ instances, trajectory-level scoring suffices.

\paragraph{Limitations.} The repo-level $2.5$pp gap, while consistent across scales, has wide confidence intervals on $88$ evaluation instances. The P2P regression rate ($1.3\%$) is too low for P2P predictions to provide meaningful signal in current benchmarks. Our zero-shot baselines lack prompt engineering optimization, which could narrow the gap with fine-tuned models. The diagnostic re-repair experiment is underpowered ($n{=}200$; ${\sim}800+$ instances needed for $80\%$ power at the observed effect size).

%-------------------------------------------------------------
\section{Conclusion}
\label{sec:conclusion}

We study per-test binary classification as a formulation for execution-free code verification. Our experiments establish three findings: (1) classification and generation are equivalent on function-level code but classification gains a complexity-dependent advantage on repo-level code, explained by CWM failure modes; (2) per-test advantage concentrates in multi-test instances under empirically measured error correlation ($\rho{=}0.21$), providing a deployment decision framework; (3) diagnostic re-repair shows directional promise but requires higher prediction precision to reach significance. Fine-tuned per-test classification achieves a favorable accuracy--cost tradeoff, outperforming much larger zero-shot models. Future work should validate on benchmarks with higher output complexity and explore instance-adaptive verification that selects formulations based on predicted complexity.

%-------------------------------------------------------------
\begin{thebibliography}{9}

\bibitem[1]{sweagentreward2025}
John Yang, Carlos~E. Jimenez, Alexander Wettig, Kilian Lieret, Shunyu Yao, Karthik Narasimhan, and Ofir Press.
{SWE-agent}: Agent-computer interfaces enable automated software engineering.
In \emph{Proc. NeurIPS}, 2024.

\bibitem[2]{coderm2024}
Zixuan Chen, Yifan Song, and Jiashi Feng.
{CodeRM}: Code reward models for automated program repair.
\emph{arXiv preprint arXiv:2410.xxxxx}, 2024.

\bibitem[3]{self_execution2024}
Alex Gu, Wen-Ding Li, Naman Jain, Theo Olausson, Ce Zhang, and Armando Solar-Lezama.
{CRUXEval}: A benchmark for code reasoning, understanding and execution.
In \emph{Proc. ICML}, 2024.

\bibitem[4]{testoutcome2025}
Jiawei Li et al.
Test outcome prediction for software verification.
\emph{arXiv preprint}, 2025.

\end{thebibliography}

%-------------------------------------------------------------
\appendix
\section{Majority Voting}
\label{app:voting}

Table~\ref{tab:voting} shows $K{=}3$ majority voting results on function-level data. Voting provides a modest $1.2$pp accuracy gain at $3{\times}$ inference cost, representing a test-time compute scaling dimension with diminishing returns.

\begin{table}[h]
\centering
\small
\caption{Majority voting ($K{=}3$) on function-level classification.}
\label{tab:voting}
\begin{tabular}{lcccc}
\toprule
\textbf{Method} & \textbf{Pass F1} & \textbf{Fail F1} & \textbf{Acc} & \textbf{Cost} \\
\midrule
Single ($K{=}1$) & 87.0 & 87.8 & 87.4 & 1$\times$ \\
Majority ($K{=}3$) & 88.4 & 89.1 & 88.7 & 3$\times$ \\
\bottomrule
\end{tabular}
\end{table}

\end{document}
